\section{Introduction}\label{sec:r31introduction}
Dynamic economic decisions require a continuation value, but the relevant
accuracy is the welfare loss of the policy induced by that continuation.
These are different objects. A small average regression loss need not control
profitable deviations at unvisited states, and an accurate policy need not
have been obtained economically. Neural Bellman Operators (NBO) separate
policy evaluation from feasible improvement so that a learned continuation
can be reused across decisions. This paper develops the corresponding
training, policy-accuracy, and work accounts within the original
controlled-economy framework.

The central question is constructive: when can a trainable continuation be
obtained with a prescribed computation budget and produce a policy whose
loss is bounded against the full admissible comparison class? We answer this
question for a square-activation neural continuation in a multisector
capital-adjustment economy, including a nonlinear recursive entropic
criterion. The result links economic primitives to a training schedule,
that schedule to a bound on hidden-weight approximation error, and that
error to full-policy accuracy. An independent verifier assesses the actual
rounded policy rather than assuming that the training argument has already
certified its implementation. The general smooth-network formulation and
the economically distinct preference and strategic applications are retained.

There are three parts to the argument. First, a centered continuation-error
bound measures precisely the part of an evaluation error that can change a
decision. Transport bounds describe changes of economic futures. These
bounds concern a candidate decision and remain valid whether its generator
is neural, conventional, or simulation based. Their policy composition must
use the implemented policy's own continuation and cover states reachable by
deviations, not only states visited by the fitted policy.

Second, an all-state comparison makes this policy link executable. In the
capital economy, controls are vectors at every one of 24 dates, Gaussian
innovations have continuous unbounded support, and future policies are
reoptimized in each counterfactual regime. Exact own-policy evaluation and
outward coefficient bounds replace an empirical state catalogue. A coercive
Bellman argument compares the returned policy with every adapted finite-cost
policy. For recursive risk preferences, a quadratic Bellman subsolution
propagates separate coefficient and constant allowances through the Gaussian
certainty equivalent. Its checked risk-domain margins replace an unspecified
global utility Lipschitz constant. Class approximation and value transfer
are zero for this directly specified discrete economy, not for an unrelated
diffusion approximation or equilibrium problem.

Third, we establish a constructive result for the trained representation.
A reference feedback and the economic primitives determine spectral envelopes,
local training tolerances, initialization scales, step sizes, and finite
iteration caps before any policy-evaluation target is computed. A joint
backward induction bounds the future policy and the target that it supplies.
The hidden-weight iteration contracts the Gram error from an isotropic
initialization, and full-action residual bounds then deliver the prescribed
recursive policy allowance. This is a statement about updated hidden
weights, not a theorem for a fixed feature dictionary relabeled as neural
training. The numerical certificate separately charges rounding and the
stated relative action-execution error.

These results distinguish three propositions that should not be conflated:
existence of an accurate returned policy, finite construction of that policy
by the specified neural algorithm, and a reduction in complete work relative
to another algorithm. The first two have explicit mathematical premises and
bounds. The third is a comparison of complete procedures at a common economic
accuracy. In a quadratic economy, conventional quadratic and structural
methods have access to the same analytic information and must remain strong
comparators. Neither a certified action nor a finite neural training bound
establishes that its neural construction is necessary or cheapest.

The evidence is organized around this distinction. The earlier finite-law
capital services isolate continuation reuse and changing futures under exact
scalar-action verification. Their adverse results are retained: adaptive
neural reuse can cost more than fresh fitting, and cached simulation can be
the least expensive accuracy-qualified procedure. A subsequent frozen design
certifies complete vector policies over the full state domain and enumerates
its declared finite training randomizer. The recursive experiment uses the
primitive-only allocation and an independent residual-centered verifier. Each
experiment keeps its own target, source identity, unsuccessful attempts,
and complete work account. Historical clocks are never replaced by newly
measured clocks, and finite-catalogue success is not presented as population
reliability.

The economic experiments are controlled computational economies, not estimates
of an observed policy reform. Technology and valuation changes in the
full-policy designs reoptimize all future controls; the earlier scalar design
conditions on specified future policies. The distinction matters both for
the economic interpretation and for the admissible comparison class.
Recursive utility, endogenous preference adjustment, temporal selves, and
games retain their original formulations and verification conditions in the
Economic Applications companion. A solved recursive capital instance does
not automatically supply equilibrium or deviation constants for the other
models.

The article presents the model, the NBO construction, the decision and policy
bounds, and the integrated evidence. The current Technical Supplement gives
the full-policy, recursive-risk, and constructive-learning proofs. Earlier
proofs, numerical designs, and all adverse tables are retained in the
Historical Technical Archive, with their original labels and source
identities. This organization preserves the scientific record without making
the sequence of repository revisions the structure of the current argument.
